\section{Problem Setup}\label{sec:setup}

\paragraph{Families and scores.}
A protein \emph{family} is a cluster of homologous wild-type sequences obtained without labels (Section~\ref{sec:data}). Family $k$ is the random object
$G_k=(D_k,N_k,(X_{k,i},Y_{k,i})_{i=1}^{N_k})$, where $X_{k,i}$ describes a missense variant, $Y_{k,i}\in\{0,1\}$ is its label ($1$ = pathogenic), and $D_k\in\mathcal{D}$ is a \emph{label-free descriptor} of the family: any function of its sequences, alignments, predictor scores, and of the fixed set of training families, but not of labels.
A predictor score is turned into $\hat p(y\mid x)$ by a base model fitted on separate training families, which we treat as fixed. The nonconformity score is $s(x,y)=1-\hat p(y\mid x)$ and $S_{k,i}=s(X_{k,i},Y_{k,i})$.
We observe calibration families $G_1,\dots,G_K$ with labels and a test family $G_{K+1}$ whose descriptor $D_{K+1}$ and variants $X_{K+1,i}$ are observed but whose labels are not.

\paragraph{Coverage targets.}
For a set-valued predictor $C:\mathcal{X}\to 2^{\{0,1\}}$, the per-family coverage of the test family is $\gamma_{K+1}(C)=N_{K+1}^{-1}\sum_{i}\1\{Y_{K+1,i}\in C(X_{K+1,i})\}$. Our primary target is \emph{family-weighted coverage}
\begin{equation}
\Cov(C)=\E\big[\gamma_{K+1}(C)\big],
\end{equation}
the probability that a uniformly random variant of a new family is covered; every family counts once, whatever its number of variants. For example, if one test family has 500 variants and nine others have 5 each, the large family carries over 90\% of the weight in variant-weighted coverage but a tenth of it in $\Cov$. Its class-conditional version for label $y$ is
$\Cov_y(C)=\E\big[N_{K+1,y}^{-1}\sum_{i:Y_{K+1,i}=y}\1\{y\in C(X_{K+1,i})\}\,\big|\,N_{K+1,y}>0\big]$, where $N_{k,y}$ counts label-$y$ variants of family $k$. For clinical use, $\Cov_1$ (pathogenic coverage) is the quantity of interest: $1-\Cov_1$ is the family-weighted \emph{pathogenic miss rate}, the rate at which a pathogenic variant's set excludes the pathogenic label (it is called benign or, rarely, receives an empty set). Variant-weighted coverage, which pools all test variants, lets families with many variants dominate; we report it but do not target it. Coverage for each individual family cannot be guaranteed without assumptions on how families differ~\cite{vovk2012conditional}, so the average over new families is the natural distribution-free target.

\paragraph{Assumptions.} We use two alternative assumptions; each result states which one it needs.
\begin{itemize}
\item[\textbf{A1}] (\emph{Family exchangeability.}) $(G_1,\dots,G_{K+1})$ is exchangeable.
\item[\textbf{A2}] (\emph{Family-level covariate shift.}) $G_1,\dots,G_K\stackrel{\text{iid}}{\sim}P$ and, independently, the test family $G_{K+1}$ and any further unlabeled target families are i.i.d.\ from $Q$, where $Q\ll P$ and $\frac{dQ}{dP}(G)=w(D)$ for a measurable $w:\mathcal{D}\to[0,\infty)$. Equivalently, the conditional law of a family's variants and labels given its descriptor is the same under $P$ and $Q$; only the distribution of descriptors shifts.
\end{itemize}
A1 holds when calibration and test families are drawn by a random split of a pool of families, regardless of dependence between families. It does not require exchangeability of variants within a family, because $\Cov$ averages over all variants of the test family (this is weaker than the hierarchical exchangeability of~\cite{lee2026hierarchical}, which targets one variant of the new group). A2 formalizes deployment on families that differ from the calibration families in label-free ways, such as being evolutionarily distant or having shallow alignments.
